% [Lynn] Methodology (rubric: 3 marks).
% Top band: features, baseline, 2+ approaches, hyperparameters AND how they were tuned, validation
%   strategy, metrics: all specified and justified. Reproducible.
% FB3: cyclical time of day; every fitted step learned from training days only; why is_weekend was dropped.
% FB4: the proposal was too dense, so use numbered steps + one table, not packed prose.
% Lec 7: baseline ladder, loss vs metric (MAE), train-vs-CV error curve (from M4).
\section{Methodology}

\subsection{Prediction setting and features}
TODO. % order-time-only rule; feature table (README F1-F9)

\subsection{Baselines and models}
TODO. All models are implemented in scikit-learn~\cite{pedregosa2011scikit}. % median predictor, platform promise; Elastic-Net, Random Forest, gradient boosting; q80 model

\subsection{Validation, tuning and metrics}
TODO. % forward-chaining CV on 17-22 Oct, test 23-24 Oct once; random search; MAE/RMSE/R^2, late rate; bootstrap CIs

% \begin{table}[t]
% \centering
% \caption{Models, tuned hyperparameters and search ranges.}
% \label{tab:models}
% \begin{tabular}{lll}
% \toprule
% Model & Hyperparameters & Search range \\
% \midrule
% ... \\
% \bottomrule
% \end{tabular}
% \end{table}
